% --- CHUNK_METADATA_START ---
% src_checksum: d5f251ce0cee0631c9c61b74a4e7807b43906d760fcd99ee52fae608e116064e
% needs_review: True
% --- CHUNK_METADATA_END ---
\chapter{Lebesgue Measure and Probability Laws}% --- CHUNK_METADATA_START ---
% src_checksum: 72866aa510b3cc88f678364c0035ffe135e58e0d2ca75a0b2e5a6cc256da9ebb
% --- CHUNK_METADATA_END ---
\sloppy% --- CHUNK_METADATA_START ---
% src_checksum: 260d0fa15a49296f8460a757791fd3eaa1b139ac0faf08bda4860e56641ee95e
% needs_review: True
% --- CHUNK_METADATA_END ---
\section{The Lebesgue measure}% --- CHUNK_METADATA_START ---
% src_checksum: 8b05d5233dd2e6363c743b62496adc4d1c86123a3c9538d497f7cbf6494ca535
% needs_review: True
% --- CHUNK_METADATA_END ---
\begin{definition}
A measure on $(\mathbb{R}, \mathcal{B}(\mathbb{R}))$ is an application $\mu$ from $\mathcal{B}(\mathbb{R})$ to $[0, +\infty[$ such that $\mu(\emptyset) = 0$ and that is $\sigma$-additive, that is:
\[
\forall (B_n)_{n \in \mathbb{N}} \in \mathcal{B}(\mathbb{R})^{\mathbb{N}}, \quad \left( \forall n \neq m, B_n \cap B_m = \emptyset \implies \mu\left(\bigcup_{n \in \mathbb{N}} B_n\right) = \sum_{n \in \mathbb{N}} \mu(B_n) \right)
\]
\end{definition}% --- CHUNK_METADATA_START ---
% src_checksum: cde0d0b21c18cb88178cea652393217764fcd0fcdf8e14ec23b7be4ad7fbe8e9
% needs_review: True
% --- CHUNK_METADATA_END ---
\begin{theorem}
There exists a unique measure $\lambda$ on $(\mathbb{R}, \mathcal{B}(\mathbb{R}))$ that satisfies:
\[
\forall a, b \in \mathbb{R}, a < b, \quad \lambda([a,b]) = b-a.
\]
This measure is called the \textbf{Lebesgue measure} and is denoted by $\lambda$.
\end{theorem}% --- CHUNK_METADATA_START ---
% src_checksum: 1f7bdfa6e57da483fc59119e303c3eebd53a061a7472a2e9b6ec0f5fd971d656
% needs_review: True
% --- CHUNK_METADATA_END ---
\begin{remark}
The existence of such a measure is by no means obvious; it is a profound mathematical result.
\end{remark}% --- CHUNK_METADATA_START ---
% src_checksum: a11b67aac444c32fe779a1dcbe54ebc2eb179e7f7f95270bfa32c836d7b6b511
% needs_review: True
% --- CHUNK_METADATA_END ---
\subsection*{Proof idea}% --- CHUNK_METADATA_START ---
% src_checksum: 37a91847a0ffd01fd6e1b27117e667c16c9b159de24574a69d27cc8629756bd7
% needs_review: True
% --- CHUNK_METADATA_END ---
For any subset $A$ of $\mathbb{R}$, we define the Lebesgue outer measure $\lambda^*(A)$ by:
\[
\lambda^*(A) = \inf \left\{ \sum_{i \in I} (b_i - a_i) ; (a_i)_{i \in I}, (b_i)_{i \in I} \text{ families of real numbers such that } \forall i \in I, a_i < b_i, A \subset \bigcup_{i \in I} ]a_i, b_i[ \right\}
\]
$\lambda^*$ is not a measure, it is subadditive but not $\sigma$-additive. We show that the restriction of $\lambda^*$ to $\mathcal{B}(\mathbb{R})$ is $\sigma$-additive and we then take $\lambda = \lambda^*|_{\mathcal{B}(\mathbb{R})}$.% --- CHUNK_METADATA_START ---
% src_checksum: 814f940a8a98d88c81473bdef6ff170f09589dbda59240f36ffb97419179c664
% needs_review: True
% --- CHUNK_METADATA_END ---
\subsection*{First properties}% --- CHUNK_METADATA_START ---
% src_checksum: 9af87430a03566e89a239ad0caa791d2ff467753d949fd9302a95bdebe4b355a
% needs_review: True
% --- CHUNK_METADATA_END ---

\begin{itemize}
    \item For all $x \in \mathbb{R}$, $\lambda(\{x\}) = 0$.
    \item If $A \subset \mathbb{R}$ is countable, then $\lambda(A) = \sum_{x \in A} \lambda(\{x\}) = 0$ by $\sigma$-additivity.
    \item There exist Borel sets $C$ in $\mathbb{R}$ that are uncountable such that $\lambda(C)=0$.
\end{itemize}
% --- CHUNK_METADATA_START ---
% src_checksum: a9eae3ce37ef3d8e70a4eef83b21b01bb79bd7e1a39f0d6aa93502bc1c1c33fc
% needs_review: True
% --- CHUNK_METADATA_END ---
\section{Uniform law on $[0,1]$}}% --- CHUNK_METADATA_START ---
% src_checksum: 39de4fd650a3e1bb91704333d23e1df10c6d548dff964df4d6a678e1dcaaf247
% needs_review: True
% --- CHUNK_METADATA_END ---
The uniform law on $[0,1]$ is the probability measure $P$ on $([0,1], \mathcal{B}([0,1]))$ obtained as the restriction of the Lebesgue measure $\lambda$ to $[0,1]$:
\[
\forall B \in \mathcal{B}([0,1]), \quad P(B) = \lambda(B).
\]
We want to construct a mathematical model of the experiment consisting of drawing a number at random in $[0,1]$, uniformly.
We take the measurable space $(\Omega, \mathcal{F}) = ([0,1], \mathcal{B}([0,1]))$. The Borel sigma-algebra $\mathcal{B}([0,1])$ is, at choice:% --- CHUNK_METADATA_START ---
% src_checksum: 0a72a7362fe9fefbfcf4580bcafb2a6dc69d2ffef7673197575d5ecefe5ad0bb
% needs_review: True
% --- CHUNK_METADATA_END ---
\begin{itemize}
    \item  the tribe generated by the sub-intervals of $[0,1]$,
\item  or the trace tribe of $\mathcal{B}(\mathbb{R})$ on $[0,1]$, that is, $\{B \cap [0,1] : B \in \mathcal{B}(\mathbb{R})\}$.
\end{itemize}% --- CHUNK_METADATA_START ---
% src_checksum: 0e37bffbee98cdef6a56dd7a88f2a8bfd61e61c6654c43cdbfc679c715074a2c
% needs_review: True
% --- CHUNK_METADATA_END ---
This space is equipped with the probability $P = \lambda|_{[0,1]}$.
We then define a random variable $X$ by setting:% --- CHUNK_METADATA_START ---
% src_checksum: a97d968cae4060870dde7816bcf56ac4e57a4fa8f0ca02358bce1b21d92f26dc
% --- CHUNK_METADATA_END ---
\begin{align*}
    X: \Omega = [0,1] &\to \mathbb{R} \\
    \omega &\mapsto X(\omega) = \omega
\end{align*}% --- CHUNK_METADATA_START ---
% src_checksum: 4713431710b2148f8a689962265428ab2ec1cad0afe63e51b69a82639ca712f3
% needs_review: True
% --- CHUNK_METADATA_END ---
It is said that $X$ is a number chosen randomly and uniformly in $[0,1]$.
The law of $X$, denoted $P_X$, is $\lambda|_{[0,1]}$. Indeed, for all $B \in \mathcal{B}([0,1])$:
\[
P_X(B) = P(X \in B) = P(\{\omega \in \Omega : X(\omega) \in B\}) = P(B) = \lambda(B).
\]% --- CHUNK_METADATA_START ---
% src_checksum: 2c695f41371b86502cdeed2274264e5e7bd54b0e9846854750374f746ef5169c
% needs_review: True
% --- CHUNK_METADATA_END ---
\section{Density Laws}% --- CHUNK_METADATA_START ---
% src_checksum: abe2c772361c59bf39588a50ff396f66dd5d2d30442fb2d09100bb34094e5db7
% needs_review: True
% --- CHUNK_METADATA_END ---
\begin{definition}
A probability law $\mu$ on $\mathbb{R}$ is said to have a \textbf{density} if there exists a positive Borel function $f$ such that:
\[
\forall A \in \mathcal{B}(\mathbb{R}), \quad \mu(A) = \int_A f d\lambda = \int_A f(x) dx.
\]
The function $f$ is called the density (or a density) of $\mu$. A density is a positive Borel function that satisfies:
\[
\int_{\mathbb{R}} f d\lambda = \int_{\mathbb{R}} f(x) dx = 1.
\]
The density of a law is not unique; it is defined up to a set of Lebesgue measure zero. Conversely, given a positive Borel function $f$ satisfying this condition, one can construct a probability law on $\mathbb{R}$ associated with $f$ using the above formula.
\end{definition}% --- CHUNK_METADATA_START ---
% src_checksum: e02106dbda30e3e9a20d16c46ff5a04b89f440b4d633d272a3da07b19e3462c3
% needs_review: True
% --- CHUNK_METADATA_END ---
\begin{example}[Exemples fondamentaux][Fundamental Examples]
\begin{itemize}
    \item \textbf{Uniform distribution} on an interval $[a,b]$, denoted $\mathcal{U}([a,b])$. It is the distribution with density $\frac{1}{b-a}\mathbb{1}_{[a,b]}(x)$.
    \item \textbf{Exponential distribution} with parameter $\alpha > 0$, denoted $\text{Exp}(\alpha)$. It is the distribution with density $\alpha e^{-\alpha x} \mathbb{1}_{\mathbb{R}_+}(x)$.
    \item \textbf{Normal distribution} or Gaussian with parameters $m, \sigma^2$, denoted $\mathcal{N}(m, \sigma^2)$. It is the distribution with density $\frac{1}{\sqrt{2\pi}\sigma} \exp\left(-\frac{(x-m)^2}{2\sigma^2}\right)$.
\end{itemize}
\end{example}% --- CHUNK_METADATA_START ---
% src_checksum: 5f98fb878187f2fcc0b3e8400ccf81d3be1f6b82c210cbd74c30b7255caa4549
% needs_review: True
% --- CHUNK_METADATA_END ---
\section{Classification of laws on $\mathbb{R}$}}% --- CHUNK_METADATA_START ---
% src_checksum: 6ee6520019f41c7119e8c1ec893d7f83836b06240f8a926eb0e01bfb4f99a6e3
% needs_review: True
% --- CHUNK_METADATA_END ---
A point $x \in \mathbb{R}$ is an \textbf{atom} of the law $\mu$ if $\mu(\{x\}) > 0$.
We say that the law $\mu$ is \textbf{diffuse} if it has no atoms.
Three "pure" types of laws are distinguished:% --- CHUNK_METADATA_START ---
% src_checksum: 084c255ce6967e351f6985b82e0df11a635a44d857b2ecd28d357ab741d068f4
% needs_review: True
% --- CHUNK_METADATA_END ---

\begin{itemize}
    \item \textbf{Discrete laws:} A measure $\mu$ is said to be discrete if there exists a finite or countable set $D$ such that $\mu(D) = 1$. We then have $\mu = \sum_{x \in D} \mu(\{x\}) \delta_x$.
    \item \textbf{Laws with density:} A measure $\mu$ is said to have a density if there exists a Borel function $f \ge 0$ such that $\forall A \in \mathcal{B}(\mathbb{R})$, $\mu(A) = \int_A f d\lambda$.
    \item \textbf{Singular laws:} A measure $\mu$ is said to be singular if it is diffuse and if there exists a set $C$ such that $\mu(C)=1$ and $\lambda(C)=0$.
\end{itemize}
% --- CHUNK_METADATA_START ---
% src_checksum: bc270db5e3f5499d33df2a7c5934f5b45c3b70e9936760762ab5022a272ea496
% needs_review: True
% --- CHUNK_METADATA_END ---
Every probability measure on $\mathbb{R}$ can be decomposed as a convex combination of these three types:
\[
\mu = \alpha_d \mu_d + \alpha_a \mu_a + \alpha_s \mu_s
\]
where $\mu_d, \mu_a, \mu_s$ are discrete, absolutely continuous, and singular, respectively, and $\alpha_d, \alpha_a, \alpha_s \in [0,1]$ with $\alpha_d + \alpha_a + \alpha_s = 1$. This representation is unique.% --- CHUNK_METADATA_START ---
% src_checksum: e99d11a8cde272e743ef8addabbf198b153668fa1236c135800dcbac52145a42
% needs_review: True
% --- CHUNK_METADATA_END ---
\section{The Lebesgue integral}% --- CHUNK_METADATA_START ---
% src_checksum: b6b71ad7f94ffe8807237c7a0272bcc7b8559ce43c7d95f28d8542434ba77d71
% needs_review: True
% --- CHUNK_METADATA_END ---
Suppose we have a theory of integration. Given a function $f \ge 0$ such that $\int_{\mathbb{R}} f(x) dx = 1$, we can define a probability measure by $\mu(A) = \int_A f(x) dx$ for all sets $A$ allowed by the theory.
For example, with the Riemann integral, we can only integrate "sufficiently regular" functions (e.g., continuous ones), and the probability is only defined on intervals or finite unions of intervals.
The preferred choice for probability theory is the Lebesgue integral, which allows integration over all Borel sets.% --- CHUNK_METADATA_START ---
% src_checksum: 74bcd1bd7ebbe65ba88b5aba23a863e9393b2a21f758bcd52d16da966b937e61
% needs_review: True
% --- CHUNK_METADATA_END ---
\subsection*{Summary of the construction of the Lebesgue integral}% --- CHUNK_METADATA_START ---
% src_checksum: 9e7f1b4645efeac5184b75776db022b8bbfe45b8cb35d08e1ac6abf8b2dcffad
% needs_review: True
% --- CHUNK_METADATA_END ---
\begin{enumerate}
    \item \textbf{Layered functions:} A Borel function $s$ is said to be layered if it takes a finite number of values. It has the form $s = \sum_{i=1}^n a_i \mathbb{1}_{A_i}$, where $a_i \in \mathbb{R}$ and $A_i \in \mathcal{B}(\mathbb{R})$. Its integral with respect to $\lambda$ is:
    \[ \int s \,d\lambda = \sum_{i=1}^n a_i \lambda(A_i). \]
    \item \textbf{Positive functions:} If $f$ is a Borel function that is positive or zero, the integral of $f$ with respect to $\lambda$ is:
    \[ \int f \,d\lambda = \sup \left\{ \int s \,d\lambda \ ; \ s \text{ layered, } 0 \le s \le f \right\}. \]
    Note that the above supremum may be infinite.
    \item \textbf{Functions of arbitrary sign:} Let $f$ be a Borel function. We say that $f$ is integrable if $\int |f| \,d\lambda < +\infty$.
    In this case, we decompose $f$ into its positive part $f^+ = \max(f,0)$ and its negative part $f^- = \max(-f,0) = -\min(f,0)$. We then define:
    \[ \int f \,d\lambda = \int f^+ \,d\lambda - \int f^- \,d\lambda. \]
    We denote $L^1(\mathbb{R}, \lambda)$ the set of $\lambda$-integrable functions.
\end{enumerate}% --- CHUNK_METADATA_START ---
% src_checksum: 281d34050f2b32573c48548acbf1ed436ef8e5ab6a29dce75865e824d5585a15
% needs_review: True
% --- CHUNK_METADATA_END ---
\section{Transfer formula}% --- CHUNK_METADATA_START ---
% src_checksum: 5f341a944e306c17050bf6bae707f4892b1410c88ba502aa9384c313749cb440
% needs_review: True
% --- CHUNK_METADATA_END ---
When working with random variables, the first step involves performing manipulations in the probability space $(\Omega, \mathcal{F}, P)$. When dealing with the specific distributions of the random variables, the computation is transferred to $\mathbb{R}$.% --- CHUNK_METADATA_START ---
% src_checksum: 147d9505d0433ee956003a02d66f8248409dc4f213c5eb4821b128e635208018
% needs_review: True
% --- CHUNK_METADATA_END ---
\begin{proposition}[Formule de transfert]
Let $(\Omega, \mathcal{F}, P)$ be a probability space, $X$ a random variable taking values in $\mathbb{R}$, and $g: \mathbb{R} \to \mathbb{R}$ a Borel function such that $g(X)$ is an integrable random variable. Then
\[ E(g(X)) = \int_{\Omega} g(X(\omega)) \,dP(\omega) = \int_{\mathbb{R}} g(x) \,dP_X(x). \]
\end{proposition}% --- CHUNK_METADATA_START ---
% src_checksum: 5ff1d5838e1ab0a36a0bda7167c1e7ad8842da14d66f341d3fbf7f617cd5b8a5
% needs_review: True
% --- CHUNK_METADATA_END ---
\begin{proof}[Idée de la preuve]
The proof proceeds in successive steps by showing that the formula holds for increasingly general classes of functions.
\begin{enumerate}
    \item \textbf{Case of an indicator function:} Let us start with $g = \mathbb{1}_B$ where $B \in \mathcal{B}(\mathbb{R})$.
    \begin{align*}
    \int_{\Omega} \mathbb{1}_B(X(\omega)) \,dP(\omega) &= P(\{\omega \in \Omega; X(\omega) \in B\}) \\
    &= P(X \in B) = P_X(B) \\
    &= \int_{\mathbb{R}} \mathbb{1}_B(x) \,dP_X(x).
    \end{align*}
    The formula is thus verified for indicator functions.
    \item \textbf{Case of a positive simple function:} By the linearity of the integral, the formula extends to positive simple functions $s = \sum a_i \mathbb{1}_{B_i}$.
    \item \textbf{Case of a positive Borel function:} We use a monotone convergence argument. Let $\mathcal{L}$ be the collection of positive Borel functions for which the transfer formula holds. We show that $\mathcal{L}$ satisfies the following properties:
    \begin{enumerate}
        \item[(i)] $\forall B \in \mathcal{B}(\mathbb{R})$, $\mathbb{1}_B \in \mathcal{L}$.
        \item[(ii)] $\forall f,g \in \mathcal{L}$, $\forall \alpha, \beta > 0$, $\alpha f + \beta g \in \mathcal{L}$ (stability under positive linear combination).
        \item[(iii)] $\mathcal{L}$ is stable under increasing limits: if $(f_n)$ is a sequence of functions in $\mathcal{L}$ such that $f_n \le f_{n+1}$ and $f_n \to f$ pointwise, then $f \in \mathcal{L}$.
    \end{enumerate}
    By the monotone class lemma, $\mathcal{L}$ contains all positive Borel functions.
    \item \textbf{General case:} For a function $g$ of any sign, we apply the result to its positive part $g^+$ and its negative part $g^-$.
\end{enumerate}
\end{proof}% --- CHUNK_METADATA_START ---
% src_checksum: af8410aed3ddb750955e0fde13d66dffc335c67520e1ba6596e4a14680c06689
% needs_review: True
% --- CHUNK_METADATA_END ---
\subsection{Application to Expectation and Variance}% --- CHUNK_METADATA_START ---
% src_checksum: b1299b232af5ffae5e2ec35dde81bf2fc7b29dbd8e77bbd0150dfe6381412033
% needs_review: True
% --- CHUNK_METADATA_END ---
Let $X$ be an integrable random variable on $(\Omega, \mathcal{F}, P)$ with values in $\mathbb{R}$.
\paragraph{Expectation}
The expectation of $X$ is defined by $E(X) = \int_{\Omega} X(\omega) dP(\omega)$. By applying the transfer formula with $g(x) = x$, we get:
\[
E(X) = \int_{\mathbb{R}} x \,dP_X(x).
\]
If the law of $X$, $P_X$, has a density $f$, then:
\[
E(X) = \int_{x \in \mathbb{R}} x f(x) dx.
\]
\paragraph{Variance}
If $X$ is square-integrable, its variance is defined by $\text{Var}(X) = E[(X-E(X))^2]$. By setting $m=E(X)$ and applying the transfer formula to the function% --- CHUNK_METADATA_START ---
% src_checksum: 03db14e41f9f859a0c9b11333fd1ab4d8664e1cbb8b2a7979d7fb5a0ee4e074f
% needs_review: True
% --- CHUNK_METADATA_END ---
$g(x) = (x-m)^2$, we obtain:
\[
\text{Var}(X) = \int_{\mathbb{R}} (x-m)^2 \,dP_X(x).
\]
If the distribution of $X$ has a density $f$, then:
\[
\sigma^2 = \text{Var}(X) = \int_{x \in \mathbb{R}} (x-m)^2 f(x) dx = \int_{x \in \mathbb{R}} x^2 f(x) dx - m^2.
\]% --- CHUNK_METADATA_START ---
% src_checksum: ccbeff0a517b2343ceba9cd049f777b90927e1319c28664e076b5c6f0bc95b75
% needs_review: True
% --- CHUNK_METADATA_END ---
\begin{example}
\begin{itemize}
    \item If $X \sim \text{Exp}(\alpha)$, then $E(X) = \frac{1}{\alpha}$ and $\text{Var}(X) = \frac{1}{\alpha^2}$.
    \item If $X \sim \mathcal{N}(m, \sigma^2)$, then $E(X) = m$ and $\text{Var}(X) = \sigma^2$.
\end{itemize}
\end{example}% --- CHUNK_METADATA_START ---
% src_checksum: 028d95f71acc039d424bb4686ee4d7eaf5fb27bac1fae1970cd155a2cfea2b23
% needs_review: True
% --- CHUNK_METADATA_END ---
\section{Dumb function method}% --- CHUNK_METADATA_START ---
% src_checksum: 6a74a786e0ed3156bc7bdc866b0e45932fa72f06116e0bd61f08f01f1ff94126
% needs_review: True
% --- CHUNK_METADATA_END ---
The basic technique for determining the law of a random variable $X$ is the mute function method. We take a bounded Borel function $f: \mathbb{R} \to \mathbb{R}$ and evaluate $E(f(X))$. The transfer formula, $E(f(X)) = \int_{\mathbb{R}} f(x) \,dP_X(x)$, allows us to identify the law $P_X$.
This method relies on the following proposition.% --- CHUNK_METADATA_START ---
% src_checksum: 47e4fe717eb3925bcb0692ad5164ccdd7a0b6f9de3ea13d536cd5f1b4142746c
% needs_review: True
% --- CHUNK_METADATA_END ---
\begin{proposition}
Let $\mu_1, \mu_2$ be two probability measures on $\mathbb{R}$ such that for any continuous and bounded function $f: \mathbb{R} \to \mathbb{R}$, we have $\int_{\mathbb{R}} f d\mu_1 = \int_{\mathbb{R}} f d\mu_2$.
Then $\mu_1 = \mu_2$.
\end{proposition}% --- CHUNK_METADATA_START ---
% src_checksum: 869b083341b170a0436976e8c3650194444d329b5f91cf7c305bfec3dc7aefd2
% needs_review: True
% --- CHUNK_METADATA_END ---
\begin{example}
Let $X \sim \mathcal{U}([0,1])$ and $k \ge 1$. What is the distribution of $Y=X^k$?
\end{example}% --- CHUNK_METADATA_START ---
% src_checksum: 24d165a86d4971e73d8742f3cc9d5adfbaa32ce206b74a236658e26e408a4b39
% needs_review: True
% --- CHUNK_METADATA_END ---
\begin{solution}
Let $f: \mathbb{R} \to \mathbb{R}$ be a continuous bounded function. Let's calculate $E(f(X^k))$.
\begin{align*}
E(f(X^k)) &= \int_{\Omega} f(X^k(\omega)) dP(\omega) && \\
           &\stackrel{\text{transfert}}{=} \int_{\mathbb{R}} f(x^k) dP_X(x) && \text{with } P_X \text{ having density } \mathbb{1}_{[0,1]}(x) \\
           &= \int_{\mathbb{R}} f(x^k) \mathbb{1}_{[0,1]}(x) dx = \int_0^1 f(x^k) dx
\end{align*}
Let's perform the change of variable $y = x^k$. Then $x = y^{1/k}$, and $dx = \frac{1}{k} y^{\frac{1}{k}-1} dy$. The integration limits remain from $0$ to $1$.
\begin{align*}
E(f(X^k)) &= \int_0^1 f(y) \cdot \frac{1}{k} y^{\frac{1}{k}-1} dy \\
           &= \int_{\mathbb{R}} f(y) \left( \frac{1}{k} y^{\frac{1}{k}-1} \mathbb{1}_{[0,1]}(y) \right) dy
\end{align*}
We recognize the form $\int_{\mathbb{R}} f(y) g(y) dy$. By the transfer method, we deduce that the distribution of $X^k$ is a density distribution, and its density is the function $g(y) = \frac{1}{k} y^{\frac{1}{k}-1} \mathbb{1}_{[0,1]}(y)$.
\end{solution}